\clearpage

\anonsection{ПРИЛОЖЕНИЕ А}
\label{app:grammar}

\subsection*{Формальная грамматика базисного подмножества Рефала-5}

Грамматика приведена в нотации EBNF. Синтаксис базисного подмножества является
точным подмножеством синтаксиса Рефала-5. В реализуемом базисном подмножестве дополнительно
поддерживается директива \prog{\$EXTERN}, нужная для согласованности с выходом
рассахаривателя \prog{r5fw-desugar}~\cite{refal5_framework_github}.

\subsection*{Лексика}

\begin{itemize}
  \item Пробелы и переводы строк разделяют токены и игнорируются.
  \item Комментарий: строка, начинающаяся с \prog{*}, до конца строки.
  \item Целое число: \prog{-?[0-9]+}.
  \item Строка: \prog{'\ldots'} либо \prog{\textquotedbl\ldots\textquotedbl} (без экранирования; допускаются любые
        символы до следующей кавычки того же вида).
  \item Идентификатор: \prog{[A-Za-z.\_\$][A-Za-z0-9.\_\$\-]*}.
  \item Зарезервированные директивы: \prog{\$ENTRY}, \prog{\$EXTERN}.
  \item Скобки и разделители: \prog{\{ \} ( ) < > = ; ,}.
\end{itemize}

\subsection*{Синтаксис (EBNF)}

\begin{lstlisting}[language={},caption={EBNF-грамматика базисного подмножества Refal-5},label={lst:ebnf}]
Program     ::= TopLevel* EOF ;

TopLevel    ::= EntryDecl | ExternDecl | FuncDef ;

EntryDecl   ::= "$ENTRY" Ident ";"
              | "$ENTRY" FuncDef ;

ExternDecl  ::= "$EXTERN" Ident ("," Ident)* ";" ;

FuncDef     ::= Ident "{" Rule+ "}" ;

Rule        ::= Pattern "=" Result ";"? ;

Pattern     ::= PatternTerm* ;
PatternTerm ::= Int | Str | QuotedSym | Var | "(" PatternTerm* ")" ;

Result      ::= ResultTerm* ;
ResultTerm  ::= Int | Str | QuotedSym | Var | Symbol
              | "(" ResultTerm* ")"
              | "<" Ident ResultTerm* ">" ;

Var         ::= ("s" | "t" | "e") "." VarName ;
VarName     ::= (Letter | Digit | "_" | "." | "$")+ ;

Symbol      ::= Ident ;
Ident       ::= Letter (Letter | Digit | "_" | "-" | "+" | ".")* ;
Int         ::= Digit+ ;
Str         ::= '"' Char* '"' ;
QuotedSym   ::= "'" Char* "'" ;
\end{lstlisting}

\subsection*{Семантика сопоставления}

\begin{itemize}
  \item \prog{s.X}~--- связывается с ровно одним \emph{атомарным} термом (символ, число или
        строка); скобочный терм не подходит.
  \item \prog{t.X}~--- связывается с ровно одним термом (включая скобочный).
  \item \prog{e.X}~--- связывается с последовательностью термов (возможно пустой).
  \item Если переменная встречается в правиле несколько раз~--- все вхождения должны быть
        связаны с одинаковыми значениями.
  \item В теле функции выбирается \emph{первое} правило, чей образец сопоставим с
        аргументами; затем результат подставляется и редуцируется.
\end{itemize}

\subsection*{Что \emph{не} поддерживается}

\begin{itemize}
  \item Условия \prog{,~e.Cond} после паттерна, блоки \prog{\{ : \}}~--- для них применяется
        внешний рассахариватель \prog{r5fw-desugar}, выход которого уже укладывается в эту
        грамматику.
  \item Экранирование в строках (\prog{\textbackslash n}, \prog{\textbackslash x42}).
  \item Макроцифры \prog{@BIG-MACRODIGIT@}.
  \item Псевдокомментарии \prog{*\$CLASSIC} / \prog{*\$EXTENDED}.
  \item Декларации \prog{*\$FROM} (директива модульной системы Рефала-5); несколько файлов
        загружаются через синтаксис \prog{+}: \prog{java -jar target/refal5q.jar a.ref+b.ref+c.ref}.
\end{itemize}

\clearpage
\anonsection{ПРИЛОЖЕНИЕ Б}
\label{app:listings}

\subsection*{Листинги ключевых функций}

В приложении приведены ключевые фрагменты исходного кода Java-реализации
с пояснительными комментариями. Каждый листинг сопровождён ссылкой на
соответствующий файл в каталоге \prog{refal5q-impl/src/main/java/refal/}.

\subsection*{Конкатенация контейнеров ExprOps.concat}

Файл \prog{refal5q-impl/src/main/java/refal/deque/ExprOps.java}.

\begin{lstlisting}[language=Java,caption={Конкатенация двух Expr с разбором всех четырёх случаев},label={lst:app-concat}]
static <T> Expr<T> concat(Expr<T> a, Expr<T> b) {
    if (a.isEmpty()) return b;
    if (b.isEmpty()) return a;
    if (a instanceof ShallowDeque<T> sa && b instanceof ShallowDeque<T> sb) return concatSS(sa, sb);
    if (a instanceof ShallowDeque<T> sa && b instanceof DeepDeque<T> db)    return concatSD(sa, db);
    if (a instanceof DeepDeque<T> da && b instanceof ShallowDeque<T> sb)    return concatDS(da, sb);
    DeepDeque<T> da = (DeepDeque<T>) a;
    DeepDeque<T> db = (DeepDeque<T>) b;
    return concatDD(da, db);
}

private static <T> Expr<T> concatDD(DeepDeque<T> a, DeepDeque<T> b) {
    int total = a.size() + b.size();
    ShallowDeque<T> ab = a.back();
    ShallowDeque<T> bf = b.front();
    int boundary = ab.size() + bf.size();

    Expr<ShallowDeque<T>> aMidPlus = a.mid();
    if (boundary == 0) {
        (*@{\color{black}\ttfamily // оба граничных буфера пусты -- ничего добавлять не нужно.}@*)
    } else if (boundary <= CAPACITY) {
        ShallowDeque<T> merged = ShallowDeque.empty();
        for (T x : iter(ab)) merged = merged.pushBackTight(x);
        for (T x : iter(bf)) merged = merged.pushBackTight(x);
        aMidPlus = aMidPlus.pushBack(merged);
    } else {
        Object[] tmp = new Object[boundary];
        int i = 0;
        for (T x : iter(ab)) tmp[i++] = x;
        for (T x : iter(bf)) tmp[i++] = x;
        int rightLen = boundary - CAPACITY;
        if (rightLen < ShallowDeque.MIN_FILL) rightLen = ShallowDeque.MIN_FILL;
        int leftLen = boundary - rightLen;
        ShallowDeque<T> m1 = ShallowDeque.ofArray(tmp, 0, leftLen);
        ShallowDeque<T> m2 = ShallowDeque.ofArray(tmp, leftLen, rightLen);
        aMidPlus = aMidPlus.pushBack(m1).pushBack(m2);
    }

    Expr<ShallowDeque<T>> newMid = concat(aMidPlus, b.mid());
    return DeepDeque.mkDeep(a.front(), newMid, b.back(), total);
}
\end{lstlisting}

Случай Deep+Deep явно показан полностью --- в нём проявляется рекурсивный
вызов \prog{concat} на средних очередях буферов, обеспечивающий
амортизированно \(O(1)\) сложность по Окасаки~\cite{okasaki_pfds}. Случаи
Shallow+Shallow, Shallow+Deep и Deep+Shallow (методы \prog{concatSS},
\prog{concatSD}, \prog{concatDS}) реализованы аналогично с расщеплением
объединённого буфера ровно пополам, если он превышает \prog{CAPACITY}.

\subsection*{Сопоставление образца с выражением Matcher.match}

Файл \prog{refal5q-impl/src/main/java/refal/runtime/Matcher.java}.

\begin{lstlisting}[language=Java,caption={Сопоставление образца с выражением (с backtracking по e-переменным)},label={lst:app-match}]
public static Optional<Bindings> match(List<Term> pattern, List<Term> expr) {
    return rec(pattern, 0, expr, 0, expr.size(), Bindings.empty());
}

private static Optional<Bindings> rec(
        List<Term> pat, int pi,
        List<Term> exp, int ei, int eEnd,
        Bindings env) {
    if (pi == pat.size()) {
        return ei == eEnd ? Optional.of(env) : Optional.empty();
    }
    Term p = pat.get(pi);

    if (p instanceof Var v && v.kind() == 'e') {
        String key = v.key();
        if (env.contains(key)) {
            (*@{\color{black}\ttfamily // Повторное вхождение: проверяем совпадение.}@*)
            List<Term> bound = env.get(key);
            int len = bound.size();
            if (ei + len > eEnd) return Optional.empty();
            for (int k = 0; k < len; k++) {
                if (!exp.get(ei + k).equals(bound.get(k))) return Optional.empty();
            }
            return rec(pat, pi + 1, exp, ei + len, eEnd, env);
        }
        (*@{\color{black}\ttfamily // Свежая e.X -- перебираем длину захвата 0, 1, ...}@*)
        for (int len = 0; ei + len <= eEnd; len++) {
            List<Term> chunk = exp.subList(ei, ei + len);
            Bindings env1 = env.with(key, List.copyOf(chunk));
            Optional<Bindings> r = rec(pat, pi + 1, exp, ei + len, eEnd, env1);
            if (r.isPresent()) return r;
        }
        return Optional.empty();
    }

    if (ei >= eEnd) return Optional.empty();
    Term x = exp.get(ei);

    if (p instanceof Var v) {
        String key = v.key();
        if (v.kind() == 's') {
            if (x instanceof Br) return Optional.empty();
            if (env.contains(key)) {
                List<Term> bound = env.get(key);
                if (bound.size() != 1 || !bound.get(0).equals(x)) return Optional.empty();
                return rec(pat, pi + 1, exp, ei + 1, eEnd, env);
            }
            return rec(pat, pi + 1, exp, ei + 1, eEnd, env.with(key, List.of(x)));
        }
        if (v.kind() == 't') {
            if (env.contains(key)) {
                List<Term> bound = env.get(key);
                if (bound.size() != 1 || !bound.get(0).equals(x)) return Optional.empty();
                return rec(pat, pi + 1, exp, ei + 1, eEnd, env);
            }
            return rec(pat, pi + 1, exp, ei + 1, eEnd, env.with(key, List.of(x)));
        }
    }

    if (p instanceof Br pBr) {
        if (!(x instanceof Br xBr)) return Optional.empty();
        Optional<Bindings> inner = rec(pBr.items(), 0, xBr.items(), 0, xBr.items().size(), env);
        if (inner.isEmpty()) return Optional.empty();
        return rec(pat, pi + 1, exp, ei + 1, eEnd, inner.get());
    }

    if (p.equals(x)) return rec(pat, pi + 1, exp, ei + 1, eEnd, env);
    return Optional.empty();
}
\end{lstlisting}

\subsection*{Подстановка результата правила Substitutor.substitute}

Файл \prog{refal5q-impl/src/main/java/refal/runtime/Substitutor.java}.

\begin{lstlisting}[language=Java,caption={Сборка результата правила через Expr.concat},label={lst:app-subst}]
public static Expr<Term> substitute(List<Term> result, Bindings env) {
    Expr<Term> out = Expr.empty();
    for (Term t : result) {
        switch (t) {
            case Var v -> {
                String key = v.key();
                if (!env.contains(key)) {
                    throw new IllegalStateException(
                            "Unbound variable in result: " + key);
                }
                out = Expr.concat(out, Expr.fromList(env.get(key)));
            }
            case Br br -> {
                Expr<Term> inner = substitute(br.items(), env);
                out = out.pushBack(new Br(inner.toList()));
            }
            case Call c -> {
                Expr<Term> argDq = substitute(c.args(), env);
                out = out.pushBack(new Call(c.fn(), argDq.toList()));
            }
            default -> out = out.pushBack(t);
        }
    }
    return out;
}
\end{lstlisting}

Использование switch-выражения с шаблонами над запечатанным
интерфейсом \prog{Term} обеспечивает компилируемую полноту разбора:
если в иерархию AST добавляется новый вид терма, компилятор Java потребует
обработать его явно.

\subsection*{Цикл нормализации Runtime.normalize и Runtime.evalCall}

Файл \prog{refal5q-impl/src/main/java/refal/runtime/Runtime.java}.

\begin{lstlisting}[language=Java,caption={Цикл нормализации (прикладной порядок, сначала крайний слева)},label={lst:app-normalize}]
public Expr<Term> normalize(Expr<Term> expr) {
    Expr<Term> out = Expr.empty();
    for (Term t : expr) {
        out = Expr.concat(out, normalizeTerm(t));
    }
    return out;
}

private Expr<Term> normalizeTerm(Term t) {
    return switch (t) {
        case Call c -> {
            bumpStep();
            Expr<Term> normalizedArgs = normalize(Expr.fromList(c.args()));
            Expr<Term> body = evalCall(c.fn(), normalizedArgs.toList());
            yield normalize(body);
        }
        case Br br -> {
            Expr<Term> inner = normalize(Expr.fromList(br.items()));
            yield Expr.singleton(new Br(inner.toList()));
        }
        default -> Expr.singleton(t);
    };
}

public Expr<Term> evalCall(String fn, List<Term> args) {
    BuiltinFn b = builtins.get(fn);
    if (b != null) {
        List<Term> r = b.apply(this, args);
        return Expr.fromList(r);
    }
    Optional<Func> f = program.getFunc(fn);
    if (f.isEmpty()) {
        throw new RefalException("Unknown function: " + fn);
    }
    for (Rule rule : f.get().rules()) {
        var bindings = Matcher.match(rule.pattern(), args);
        if (bindings.isPresent()) {
            return Substitutor.substitute(rule.result(), bindings.get());
        }
    }
    throw new RefalException("No matching rule for " + fn + " on args=" + args);
}
\end{lstlisting}

\clearpage
\anonsection{ПРИЛОЖЕНИЕ В}
\label{app:reproduce}

\subsection*{Команды воспроизведения тестов и сравнительного прогона}

Все команды выполняются из каталога \prog{refal5q-impl/}. Предполагаются
установленными Java~25~LTS~\cite{java_se25_spec} (с
переменной \prog{JAVA\_HOME}, указывающей на JDK) и Apache Maven~3.9 или
выше~\cite{maven_guide}.

\subsection*{Установка инструментов}

\begin{lstlisting}[language=bash,caption={Проверка инструментов перед сборкой}]
$ java -version       (*@{\color{black}\ttfamily \# ожидается 25.x}@*)
$ javac -version      (*@{\color{black}\ttfamily \# ожидается 25.x}@*)
$ mvn -version        (*@{\color{black}\ttfamily \# ожидается 3.9.x или новее}@*)
\end{lstlisting}

Зависимости проекта (JUnit~5, AssertJ, JaCoCo, Checkstyle) подтягиваются
Maven автоматически из локального репозитория \prog{\textasciitilde/.m2/}; явная
установка не требуется.

\subsection*{Полная сборка с тестами и проверками}

\begin{lstlisting}[language=bash,caption={Полная сборка проекта}]
$ mvn clean verify
\end{lstlisting}

Команда выполняет:
\begin{itemize}
  \item компиляцию (\prog{maven-compiler-plugin}, целевой релиз --- 25);
  \item запуск всех \(143\) тестов JUnit~5 (\prog{maven-surefire-plugin});
  \item построение самодостаточного JAR-файла \prog{target/refal5q.jar} (\prog{maven-shade-plugin});
  \item проверку покрытия (\prog{jacoco-maven-plugin}; пороги
        \(\geq 70\,\%\) по инструкциям и \(\geq 55\,\%\) по веткам для
        пакетов \prog{refal.deque} и \prog{refal.runtime});
  \item проверку стиля кода (\prog{maven-checkstyle-plugin}, конфигурация
        \prog{checkstyle.xml}).
\end{itemize}

При успешном завершении в каталоге \prog{target/site/jacoco/} лежит HTML-отчёт о
покрытии, в \prog{target/} --- jar-артефакт.

\subsection*{Запуск программы}

\begin{lstlisting}[language=bash,caption={Запуск программы из файла}]
$ java -jar target/refal5q.jar bench/hello.ref --expr "'World'"
\end{lstlisting}

При наличии \prog{\$ENTRY} в файле опция \prog{--entry} необязательна.
Опция \prog{--step-limit~N} устанавливает лимит шагов редукции
(по умолчанию \(100\,000\)).

\subsection*{Workflow с внешним рассахаривателем}

Если исходный файл написан с использованием расширенного синтаксиса
Рефала-5 (условия \prog{,~e.Cond}, блоки \prog{\{ : \}}), то перед запуском
его пропускают через рассахариватель из проекта
\prog{refal-5-framework}~\cite{refal5_framework_github}:

\begin{lstlisting}[language=bash,caption={Полный workflow с рассахаривателем}]
$ refgo r5fw-desugar program.ref program-desugared.ref
$ java -jar target/refal5q.jar program-desugared.ref --entry Go
\end{lstlisting}

\subsection*{Сравнительные прогоны трёх реализаций Рефала-5}

Сравнительные прогоны из подраздела~\ref{subsec:benchmarks} автоматизированы
двумя PowerShell-скриптами. Оба ожидают установленный Refal-5~PZ~\cite{refal5_home}
(дистрибутив Version~PZ с сайта ИПМ им.~А.~П.~Ершова; пути к
\prog{refgo.exe} и \prog{refc.exe} задаются в начале скрипта) и опционально~---
Refal-5J при включённой опции \prog{-WithR5J}.

Скрипт \prog{bench/loop-select-bench.ps1} прогоняет синтетический тест
Loop/Select при значениях $N$ от 100 до 20\,000:

\begin{lstlisting}[language=bash,caption={Запуск бенчмарка Loop/Select}]
$ powershell -File bench\loop-select-bench.ps1 -WithPZ -WithR5J
\end{lstlisting}

Скрипт \prog{bench/desugar-bench.ps1} измеряет время рассахаривания
\prog{R5FW-Parser.ref}:

\begin{lstlisting}[language=bash,caption={Запуск бенчмарка рассахаривателя}]
$ powershell -File bench\desugar-bench.ps1 -WithPZ -WithR5J
\end{lstlisting}

По умолчанию выполняется \(1\) прогревочный запуск и \(3\) измеряемых, в
отчёт идёт медиана; параметры регулируются опциями \prog{-WarmupRuns} и
\prog{-MeasuredRuns}. Результаты Loop/Select сохраняются в
\prog{bench/loop-select-results.csv}; график для основного текста строится
отдельно командой \prog{python bench/plot-benchmark.py --out bench/loop-select-graph.pdf}.
Время рассахаривания записывается в \prog{bench/desugar-results.csv}.

\clearpage
\anonsection{ПРИЛОЖЕНИЕ Г}
\label{app:loopselect-src}

\subsection*{Исходный код теста Loop/Select}

Тестовая программа сравнительного прогона из подраздела~\ref{subsec:benchmarks}.
Файл \prog{refal5q-impl/bench/loop-select.ref}. Аргумент \prog{<Arg 1>}~---
число итераций $N$; на каждой итерации функция \prog{Loop} строит новый
скобочный терм из растущего аккумулятора \prog{e.Acc}, а \prog{Select}
выбирает одну из двух альтернатив в зависимости от направления.

\begin{lstlisting}[language=Refal,caption={Тестовая программа Loop/Select},label={lst:app-loopselect}]
$ENTRY Go {
  = <TimeElapsed 0>
    <Loop Left () <Numb <Arg 1>>>
    <Prout <TimeElapsed>>;
}

Loop {
  s.Side (e.Acc) 0 = /* done */;
  s.Side (e.Acc) s.Rest
    = <Loop <Select s.Side (e.Acc A) (e.Acc B)> <- s.Rest 1>>;
}

Select {
  Left  t.Left t.Right = Right t.Left;
  Right t.Left t.Right = Left  t.Right;
}
\end{lstlisting}
